For (b), if \(x''\in A''{}^G(k)\), its inverse image in \(A\)
is a trivial principal homogeneous space under \(A'\) (because
\(A(k)\to A''(k)\) is assumed to be surjective) on which \(G\)
acts, and therefore defines an element
\(d(x'')\in H^1(G,A')\).  The exactness of the sequence (1) is
then immediate.

\SGANextSectionTarget{sga3:xvii:subsection:a-2}\SGASubsection{A.2}{The Group
\texorpdfstring{\(\operatorname{Ext}_{\mathrm{alg}}(G,A)\)}
{Ext alg (G,A)}}
\SGARefTarget{sga3:xvii:appendix:A:subsection:2}

Let \(A\) and \(G\) be two algebraic \(k\)-groups, and let \(E\)
and \(E'\) be two algebraic \(k\)-groups that are extensions of
\(G\) by \(A\).  These two extensions are said to be
\emph{isomorphic} if there is a \(k\)-morphism of groups
\(u:E\to E'\) that makes the following diagram commutative:
\SGARefTarget{sga3:xvii:appendix:A:diagram:A:2:isomorphic-extensions}
\[
\begin{tikzcd}[column sep=large,row sep=large]
  {} & E \arrow[dr] \arrow[d,"u"] & {} \\
  A \arrow[ur] \arrow[dr] & {} & G \\
  {} & E' \arrow[ur] & {}
\end{tikzcd}
\qquad .
\]

The group \(E\) acts on \(A\) by inner automorphisms.  If \(A\)
is commutative, this action factors through \(G\), so that \(G\)
acts on \(A\).  Conversely, if
\(A\in\operatorname{Ob}(\mathrm{Ab})_G\), we denote by
\(\operatorname{Ext}_{\mathrm{alg}}(G,A)\) the set of classes of
algebraic extensions \(E\) of \(G\) by \(A\) for which the
action of \(G\) on \(A\) defined by \(E\) agrees with the action
coming from the structure of \(A\) as an object of
\((\mathrm{Ab})_G\).

The set \(\operatorname{Ext}_{\mathrm{alg}}(G,A)\) is naturally a
bifunctor, covariant in \(A\) and contravariant in \(G\).  More
precisely:
\begin{enumerate}
\item[(a)] \SGARefTarget{sga3:xvii:appendix:A:subsection:2:item:a}
  if \(f:A\to B\) is a morphism in \((\mathrm{Ab})_G\), and if
  \(E\) represents an element of
  \(\operatorname{Ext}_{\mathrm{alg}}(G,A)\), then
  \(f_*(E)\in\operatorname{Ext}_{\mathrm{alg}}(G,B)\) is defined
  to be the class of the extension of \(G\) by \(B\) equal to
  the quotient of the semidirect product \(B\mathbin{\cdot}E\)
  (where \(E\) acts on \(B\) through \(G\)) by the algebraic
  subgroup that is the image of \(A\) under the morphism
  \((f,i)\).  This quotient is representable by Exposé VI\textsubscript{A},
  \S5, and there is therefore a commutative diagram
  \SGARefTarget{sga3:xvii:appendix:A:diagram:A:2:pushout}
  \[
  \begin{tikzcd}[column sep=large,row sep=large]
    1 \arrow[r] &
    A \arrow[r,"i"] \arrow[d,"f"'] &
    E \arrow[r] \arrow[d] &
    G \arrow[r] \arrow[d,"1_G"] &
    1 \\
    1 \arrow[r] &
    B \arrow[r] &
    f_*(E) \arrow[r] &
    G \arrow[r] &
    1
  \end{tikzcd}
  \qquad .
  \]

\item[(b)] \SGARefTarget{sga3:xvii:appendix:A:subsection:2:item:b}
  if \(g:H\to G\) is a \(k\)-morphism of algebraic
  \(k\)-groups, and if \(E\) is an extension of \(G\) by \(A\),
  the fiber product \(E\times_G H\) is naturally an extension
  of \(H\) by \(A\), denoted by \(g^*(E)\).  Thus there is a
  commutative diagram
  \SGARefTarget{sga3:xvii:appendix:A:diagram:A:2:pullback}
  \[
  \begin{tikzcd}[column sep=large,row sep=large]
    1 \arrow[r] &
    A \arrow[r] \arrow[d,"1_A"'] &
    g^*(E) \arrow[r] \arrow[d] &
    H \arrow[r] \arrow[d,"g"] &
    1 \\
    1 \arrow[r] &
    A \arrow[r] &
    E \arrow[r] &
    G \arrow[r] &
    1
  \end{tikzcd}
  \qquad .
  \]
\end{enumerate}

By adapting the proofs given in J.-P. Serre,
\emph{Groupes algébriques et corps de classe}, Chapter VII, one
endows \(\operatorname{Ext}_{\mathrm{alg}}(G,A)\) with a natural
abelian-group structure, functorial in \(A\) and \(G\).
